% ============================================================
% SECTION 3 — Methodology
% ============================================================

\section{Methodology}
\label{sec:methodology}

We constrain the CCR-modified primordial power spectrum against
Planck 2018 data using Markov Chain Monte Carlo (MCMC) sampling.

\subsection{Modified primordial spectrum implementation}

The CCR-modified primordial scalar power spectrum, eq.~(\ref{eq:P_CCR}),
is implemented as an external primordial power spectrum provider within
the Cobaya sampling framework~\cite{Torrado:2021}. For each point in
parameter space, we evaluate $\mathcal{P}_{\rm CCR}(k)$ on a
logarithmically spaced grid of 1500 points over
$k \in [10^{-6},\,3.2]\;\mathrm{Mpc}^{-1}$, where the comoving cutoff
$k_c = k_c(D)$ is computed from eq.~(\ref{eq:kc_comoving}) for each
sampled value of $\ln D$.

The modified spectrum is passed to CAMB~1.6.5~\cite{Lewis:1999} via
the \texttt{SplinedInitialPower} interface, which performs cubic spline
interpolation in $\log k$. CAMB then computes the transfer functions
and angular power spectra $C_\ell^{TT}$, $C_\ell^{TE}$, and
$C_\ell^{EE}$ up to $\ell_{\max} = 2500$ with lensing potential
accuracy set to unity. We verify that in the limit $D \to \infty$ (equivalently $k_c \to 0$),
the standard $\Lambda$CDM power spectrum is recovered to a relative 
precision of $\lesssim 2 \times 10^{-4}$ across all multipoles.


\subsection{Planck 2018 likelihoods}

We employ the following Planck 2018
likelihoods~\cite{Planck:2018likelihood}:
\begin{itemize}
  \item \textbf{Low-$\ell$ TT} ($\ell = 2$--$29$): the Commander
    component-separation likelihood, which provides the primary
    sensitivity to the CCR suppression signal.
  \item \textbf{Low-$\ell$ EE} ($\ell = 2$--$29$): the SimAll
    likelihood, which constrains the optical depth~$\tau$.
  \item \textbf{High-$\ell$ TT+TE+EE} ($\ell = 30$--$2508$): the
    Plik lite nuisance-marginalized likelihood, which constrains the
    standard $\Lambda$CDM parameters and ensures the CCR modification
    does not degrade the fit at small angular scales.
\end{itemize}
The combination of low-$\ell$ and high-$\ell$ likelihoods is essential:
the CCR signal manifests exclusively at $\ell \lesssim 30$, while the
high-$\ell$ data anchor the cosmological parameters.

\subsection{Parameter space and priors}

The full parameter space comprises eight dimensions: six standard
$\Lambda$CDM parameters and two CCR parameters. The priors are
summarised in table~\ref{tab:priors}.

\begin{table}[t]
\centering
\begin{tabular}{lcc}
\toprule
Parameter & Prior & Description \\
\midrule
$\log_{10}(\ln D)$ & $\mathcal{U}[12,\,16]$ & Hilbert space dimension \\
$\alpha$ & $\mathcal{U}[1,\,6]$ & Suppression shape \\
$\Omega_b h^2$ & $\mathcal{U}[0.005,\,0.1]$ & Baryon density \\
$\Omega_c h^2$ & $\mathcal{U}[0.001,\,0.99]$ & CDM density \\
$H_0$ & $\mathcal{U}[40,\,100]$ & Hubble constant \\
$\tau$ & $\mathcal{N}(0.0544,\,0.0073)$ & Optical depth \\
$A_s$ & $\mathcal{U}[5 \times 10^{-10},\,5 \times 10^{-9}]$ & Scalar amplitude \\
$n_s$ & $\mathcal{U}[0.8,\,1.2]$ & Spectral index \\
\bottomrule
\end{tabular}
\caption{Prior distributions for the MCMC analysis.
  $\mathcal{U}[a,b]$ denotes a uniform prior on $[a,b]$ and
  $\mathcal{N}(\mu,\sigma)$ a Gaussian prior with mean $\mu$ and
  standard deviation $\sigma$. The Hilbert space dimension is sampled
  in $\log_{10}(\ln D)$ to efficiently explore the large dynamic range.
  We additionally fix $\sum m_\nu = 0.06\;\mathrm{eV}$ and
  $\Omega_k = 0$.}
\label{tab:priors}
\end{table}

The prior on $\log_{10}(\ln D)$ spans four orders of magnitude, from
$\ln D = 10^{12}$ (strong suppression,
$k_c \sim 7 \times 10^{-3}\;\mathrm{Mpc}^{-1}$) to $\ln D = 10^{16}$
(negligible suppression, effectively $\Lambda$CDM). The prior on
$\alpha$ encompasses the range from gradual ($\alpha = 1$, exponential)
to sharp ($\alpha = 6$, step-like) suppression profiles.

\subsection{MCMC sampling and convergence}

We sample the posterior distribution using the Metropolis--Hastings
algorithm as implemented in Cobaya~3.6.1~\cite{Torrado:2021}, with the
fast-dragging scheme enabled to exploit the hierarchy of evaluation
speeds. The CCR primordial spectrum evaluation
($\sim 9 \times 10^3$ evaluations per second) is significantly faster
than the CAMB transfer function computation ($\sim 5$ evaluations per
second), allowing the fast parameters ($\log_{10}(\ln D)$, $\alpha$,
$A_s$, $A_{\rm planck}$) to be oversampled by a factor of eight relative
to the slow cosmological parameters ($\Omega_b h^2$, $\Omega_c h^2$,
$H_0$, $\tau$, $n_s$).

The proposal covariance matrix is initialised from the Planck 2018 base
covariance for the $\Lambda$CDM parameters and learned adaptively for
the CCR parameters during a burn-in phase of 300 accepted steps.
Convergence is assessed via the Gelman--Rubin statistic $R - 1$, with
a target of $R - 1 < 0.01$. The final chains achieve $R - 1 < 0.005$ after approximately $24\,500$ accepted samples (total weight $82\,649$), with a sustained acceptance rate of $\sim 30$\%.

\subsection{Model comparison}

To assess whether the CCR modification is statistically justified, we
employ three complementary criteria:
\begin{enumerate}
  \item The \textbf{Savage--Dickey density ratio}~\cite{Dickey:1971},
    which estimates the Bayes factor by comparing the marginal posterior
    density at the $\Lambda$CDM limit ($\log_{10}(\ln D) \to 16$) to
    the prior density.
  \item The \textbf{Akaike Information Criterion} (AIC), defined as
    $\mathrm{AIC} = \chi^2_{\min} + 2k$, where $k$ is the number of
    free parameters.
  \item The \textbf{Bayesian Information Criterion} (BIC), defined as
    $\mathrm{BIC} = \chi^2_{\min} + k \ln N_{\rm data}$, where
    $N_{\rm data} \approx 669$ is the effective number of data points.
\end{enumerate}
The $\Lambda$CDM model has $k = 6$ parameters, while the CCR model has
$k = 8$ (adding $\log_{10}(\ln D)$ and~$\alpha$). These criteria
penalise model complexity and allow a quantitative assessment of whether
the two additional CCR parameters are warranted by the data.
